% TOP_PROBLEM: {"id": 10000033, "title": "No critical infinite cluster on transitive graphs", "category_id": 19, "category": {"id": 19, "name": "probability", "display_name": "Probability", "description": "Problems involving probability theory, stochastic processes, and random structures.", "slug": "probability", "order_index": 19, "created_at": "2026-07-31T15:26:25.671Z"}, "status": "open"}
% FIRST_POSTED: 2026-10-06
% ATTEMPT_STATUS: unresolved
% OPENAI_CLAIM: {"schema_version":1,"source_repo":"https://github.com/openai/math","source_commit":"adc7f1241b42e322a6451854ab7e4b4c146bf78a","release_date":"2026-10-06","match":"partial","resolution":"partial","independently_reviewed":false,"second_pass":"n/a","families":[{"family":"213","title":"Critical percolation on every quasi-transitive graph","subject":"Probability and statistical mechanics","manuscripts":[{"title":"Critical bond and site percolation on the cubic lattice","pdf_path":"preprints/Critical-bond-and-site-percolation-on-the-cubic-lattice-September-24-2026/paper.pdf","readme_path":"preprints/Critical-bond-and-site-percolation-on-the-cubic-lattice-September-24-2026/README.md","date":"2026-09-24","theorem_ref":"Theorem (label thm:main), preprints/Critical-bond-and-site-percolation-on-the-cubic-lattice-September-24-2026/build/introduction.tex:17"},{"title":"No percolation at criticality on quasi-transitive graphs","pdf_path":"preprints/No-percolation-at-criticality-on-quasi-transitive-graphs-September-24-2026/paper.pdf","readme_path":"preprints/No-percolation-at-criticality-on-quasi-transitive-graphs-September-24-2026/README.md","date":"2026-09-24","theorem_ref":"Theorem (label thm:main), preprints/No-percolation-at-criticality-on-quasi-transitive-graphs-September-24-2026/build/sections/introduction.tex:26"}],"lean":{"doc_path":"lean/docs/213.md","comparators":["lean/ComparatorChallenges/CriticalPercolation.json","lean/ComparatorChallenges/CriticalPercolation.lean","lean/ComparatorChallenges/CriticalZ3.json","lean/ComparatorChallenges/CriticalZ3.lean"],"declarations":["OAI.CriticalZ3.critical_no_infinite"],"covers_main_theorem":true},"reasoning_trace":null},{"family":"214","title":"The Benjamini–Schramm nonuniqueness conjecture","subject":"Probability and statistical mechanics","manuscripts":[{"title":"Nonuniqueness of percolation on nonamenable quasi-transitive graphs","pdf_path":"preprints/Nonuniqueness-of-percolation-on-nonamenable-quasi-transitive-graphs-September-24-2026/paper.pdf","readme_path":"preprints/Nonuniqueness-of-percolation-on-nonamenable-quasi-transitive-graphs-September-24-2026/README.md","date":"2026-09-24","theorem_ref":"Theorem (label thm:main), preprints/Nonuniqueness-of-percolation-on-nonamenable-quasi-transitive-graphs-September-24-2026/build/sections/introduction.tex:58"}],"lean":{"doc_path":"lean/docs/214.md","comparators":["lean/ComparatorChallenges/BenjaminiSchramm.json","lean/ComparatorChallenges/BenjaminiSchramm.lean","lean/ComparatorChallenges/CayleyPercolation.json","lean/ComparatorChallenges/CayleyPercolation.lean"],"declarations":[],"covers_main_theorem":true},"reasoning_trace":null}]}
\documentclass[11pt]{article}
\usepackage[margin=25mm]{geometry}
\usepackage{fontspec}
\setmainfont{Latin Modern Roman}
\usepackage{amsmath,amssymb,mathtools}
\usepackage{unicode-math}
\setmathfont{Latin Modern Math}
\usepackage{enumitem}
\usepackage{xurl,hyperref,bookmark}
\hypersetup{colorlinks=true,urlcolor=blue}
\setcounter{secnumdepth}{0}
\setlength{\parindent}{0pt}
\setlength{\parskip}{5pt}
\setlength{\emergencystretch}{3em}
\begin{document}
\section*{No critical infinite cluster on transitive graphs}

\subsection{OpenAI's claim}
The linked manuscripts present the following claims by OpenAI: Both bond and site percolation on $\mathbb Z^3$ are claimed to have only finite clusters at their respective critical probabilities; Critical bond percolation is claimed to have no infinite component on every infinite locally finite quasi-transitive graph with $p_c<1$; Every nonamenable locally finite quasi-transitive graph is claimed to satisfy $p_c<p_{2\to2}\leq p_u$, yielding a nonempty phase with infinitely many infinite clusters. Together they provide the following support within the scopes described below.

\subsection{Outline of the argument}
Independent hyperedges provide a finite connection inequality that accommodates the two models. Finite-scale extension bounds allow an adaptive exploration to maintain successful connections; assuming critical percolation would make this construction survive below the threshold, a contradiction.

The proof separates growth regimes and uses finitary structural reduction in the polynomial-growth case. Two coordinate moves are glued through fresh corridors using a conditional connection inequality; adaptive exploration would then produce subcritical percolation if an infinite critical cluster existed.

For the new unimodular case, local endpoint maps control moments of finite critical clusters. A deformation lowers susceptibility while retaining a large connectivity-operator norm, thereby proving the strict operator-threshold separation; the nonunimodular case uses the established operator result.

\subsection{Scope relative to this problem}
The theorem proves critical nonpercolation for nearest-neighbor bond and site percolation on Z3 only. The problem statement quantifies every infinite vertex-transitive graph, which is a much larger class. The theorem covers infinite connected locally finite quasi-transitive graphs only when pc<1. The literal problem statement does not exclude pc=1 or non-locally-finite graphs; Z has pc=1 and percolates at pc. The finite critical L2 two-point norm on nonamenable locally finite quasi-transitive graphs implies critical nonpercolation in that class, but does not prove the problem statement assertion for all infinite transitive graphs, especially its unexcluded pc=1 case. This entry records partial support; complete coverage of the question is not established by this catalogue.

\subsection{Formal verification}
For family 213, the Lean scope covers cubic-lattice bond and site critical nonpercolation, and quasi-transitive bond graphs when $p_c<1$. The listed Comparator files are \texttt{CriticalPercolation.lean}, \texttt{CriticalZ3.lean}. Named declarations include \texttt{OAI.CriticalZ3.critical\_no\_infinite}. For family 214, the documented statements include the strict operator-threshold separation, the nonunique percolation interval, and selected critical moment and tail bounds, with a Cayley-graph specialization. The listed Comparator files are \texttt{BenjaminiSchramm.lean}, \texttt{CayleyPercolation.lean}. The inventory does not supply a named declaration for these Comparator entries. This site has not executed the Lean code. The scope descriptions do not certify the remaining manuscript arguments.

\subsection{Status caveat}
OpenAI attributes these manuscripts to work by an internal model. Its release cautions that arguments without formalization can contain mistakes. This catalogue has not independently checked the proofs; the labels report the mathematical claims attributed to OpenAI.

\subsection{Sources}
\begin{itemize}
\item OpenAI, \emph{Critical bond and site percolation on the cubic lattice}, 2026-09-24; citation key \texttt{OAI:Critical-bond-and-site-percolation-on-the-cubic-lattice-September-24-2026}. \href{https://github.com/openai/math/blob/adc7f1241b42e322a6451854ab7e4b4c146bf78a/preprints/Critical-bond-and-site-percolation-on-the-cubic-lattice-September-24-2026/paper.pdf}{PDF (GitHub)}; \href{https://github.com/openai/math/raw/adc7f1241b42e322a6451854ab7e4b4c146bf78a/preprints/Critical-bond-and-site-percolation-on-the-cubic-lattice-September-24-2026/paper.pdf}{PDF (direct)}; \href{https://github.com/openai/math/blob/adc7f1241b42e322a6451854ab7e4b4c146bf78a/preprints/Critical-bond-and-site-percolation-on-the-cubic-lattice-September-24-2026/README.md}{Manuscript page and citation}; \href{https://github.com/openai/math/blob/main/preprints/Critical-bond-and-site-percolation-on-the-cubic-lattice-September-24-2026/paper.pdf}{Latest version}.
\item OpenAI, \emph{No percolation at criticality on quasi-transitive graphs}, 2026-09-24; citation key \texttt{OAI:No-percolation-at-criticality-on-quasi-transitive-graphs-September-24-2026}. \href{https://github.com/openai/math/blob/adc7f1241b42e322a6451854ab7e4b4c146bf78a/preprints/No-percolation-at-criticality-on-quasi-transitive-graphs-September-24-2026/paper.pdf}{PDF (GitHub)}; \href{https://github.com/openai/math/raw/adc7f1241b42e322a6451854ab7e4b4c146bf78a/preprints/No-percolation-at-criticality-on-quasi-transitive-graphs-September-24-2026/paper.pdf}{PDF (direct)}; \href{https://github.com/openai/math/blob/adc7f1241b42e322a6451854ab7e4b4c146bf78a/preprints/No-percolation-at-criticality-on-quasi-transitive-graphs-September-24-2026/README.md}{Manuscript page and citation}; \href{https://github.com/openai/math/blob/main/preprints/No-percolation-at-criticality-on-quasi-transitive-graphs-September-24-2026/paper.pdf}{Latest version}.
\item Family 213: \href{https://github.com/openai/math/blob/adc7f1241b42e322a6451854ab7e4b4c146bf78a/lean/docs/213.md}{Lean scope}.
\item \href{https://github.com/openai/math/blob/adc7f1241b42e322a6451854ab7e4b4c146bf78a/lean/ComparatorChallenges/CriticalPercolation.json}{Comparator CriticalPercolation.json}.
\item \href{https://github.com/openai/math/blob/adc7f1241b42e322a6451854ab7e4b4c146bf78a/lean/ComparatorChallenges/CriticalPercolation.lean}{Comparator CriticalPercolation.lean}.
\item \href{https://github.com/openai/math/blob/adc7f1241b42e322a6451854ab7e4b4c146bf78a/lean/ComparatorChallenges/CriticalZ3.json}{Comparator CriticalZ3.json}.
\item \href{https://github.com/openai/math/blob/adc7f1241b42e322a6451854ab7e4b4c146bf78a/lean/ComparatorChallenges/CriticalZ3.lean}{Comparator CriticalZ3.lean}.
\item OpenAI, \emph{Nonuniqueness of percolation on nonamenable quasi-transitive graphs}, 2026-09-24; citation key \texttt{OAI:Nonuniqueness-of-percolation-on-nonamenable-quasi-transitive-graphs-September-24-2026}. \href{https://github.com/openai/math/blob/adc7f1241b42e322a6451854ab7e4b4c146bf78a/preprints/Nonuniqueness-of-percolation-on-nonamenable-quasi-transitive-graphs-September-24-2026/paper.pdf}{PDF (GitHub)}; \href{https://github.com/openai/math/raw/adc7f1241b42e322a6451854ab7e4b4c146bf78a/preprints/Nonuniqueness-of-percolation-on-nonamenable-quasi-transitive-graphs-September-24-2026/paper.pdf}{PDF (direct)}; \href{https://github.com/openai/math/blob/adc7f1241b42e322a6451854ab7e4b4c146bf78a/preprints/Nonuniqueness-of-percolation-on-nonamenable-quasi-transitive-graphs-September-24-2026/README.md}{Manuscript page and citation}; \href{https://github.com/openai/math/blob/main/preprints/Nonuniqueness-of-percolation-on-nonamenable-quasi-transitive-graphs-September-24-2026/paper.pdf}{Latest version}.
\item Family 214: \href{https://github.com/openai/math/blob/adc7f1241b42e322a6451854ab7e4b4c146bf78a/lean/docs/214.md}{Lean scope}.
\item \href{https://github.com/openai/math/blob/adc7f1241b42e322a6451854ab7e4b4c146bf78a/lean/ComparatorChallenges/BenjaminiSchramm.json}{Comparator BenjaminiSchramm.json}.
\item \href{https://github.com/openai/math/blob/adc7f1241b42e322a6451854ab7e4b4c146bf78a/lean/ComparatorChallenges/BenjaminiSchramm.lean}{Comparator BenjaminiSchramm.lean}.
\item \href{https://github.com/openai/math/blob/adc7f1241b42e322a6451854ab7e4b4c146bf78a/lean/ComparatorChallenges/CayleyPercolation.json}{Comparator CayleyPercolation.json}.
\item \href{https://github.com/openai/math/blob/adc7f1241b42e322a6451854ab7e4b4c146bf78a/lean/ComparatorChallenges/CayleyPercolation.lean}{Comparator CayleyPercolation.lean}.
\item \href{https://github.com/openai/math/blob/adc7f1241b42e322a6451854ab7e4b4c146bf78a/CONTENTS.md}{OpenAI catalogue entry}.
\end{itemize}
\end{document}
